\section{列表}
\label{sec:00-lists}

\subsection{无序列表与嵌套}

写法：

\begin{lstlisting}[language=TeX]
\begin{itemize}
  \item 第一点
  \item 第二点
    \begin{itemize}
      \item 需要展开时再嵌一层
      \item 一般不要超过两层
    \end{itemize}
\end{itemize}
\end{lstlisting}

效果：

\begin{itemize}
  \item 第一点
  \item 第二点
    \begin{itemize}
      \item 需要展开时再嵌一层
      \item 一般不要超过两层
    \end{itemize}
\end{itemize}

\subsection{有序列表}

\begin{lstlisting}[language=TeX]
\begin{enumerate}
  \item 先装 TeX Live
  \item 再编译：\verb|scripts/build.sh|
  \item 最后检查 \texttt{build/main.log}
\end{enumerate}
\end{lstlisting}

效果：

\begin{enumerate}
  \item 先装 TeX Live
  \item 再编译：\verb|scripts/build.sh|
  \item 最后检查 \texttt{build/main.log}
\end{enumerate}

需要从指定编号开始时，写成 \verb|\begin{enumerate}[start=3]|。

\subsection{术语对照列表}

定义型的对照用 \texttt{description}，术语加粗放在方括号里：

\begin{lstlisting}[language=TeX]
\begin{description}
  \item[ISA] 指令集架构，软件与硬件之间的契约。
  \item[SoC] 片上系统，把处理器核、外设、存储器接口集成在一颗芯片上。
\end{description}
\end{lstlisting}

效果：

\begin{description}
  \item[ISA] 指令集架构，软件与硬件之间的契约。
  \item[SoC] 片上系统，把处理器核、外设、存储器接口集成在一颗芯片上。
\end{description}

\subsection{实验验收清单}

每节的实验都建议给一份可以逐项打勾的清单，用数学符号做方框：

\begin{lstlisting}[language=TeX]
\begin{itemize}
  \item[$\square$] 编译通过，没有错误
  \item[$\boxtimes$] 波形符合预期
\end{itemize}
\end{lstlisting}

效果：

\begin{itemize}
  \item[$\square$] 编译通过，没有错误
  \item[$\boxtimes$] 波形符合预期
\end{itemize}
